\chapter*{Tesis, método y dirección de la explicación}
\addcontentsline{toc}{chapter}{Tesis, método y dirección de la explicación}

\section*{La dirección causal}

La exposición parte de HMT\@. La matemática clásica proporciona el lenguaje de
demostración; la historia de la matemática muestra dónde surgieron los
problemas; la física conocida suministra controles, representaciones y
reconocimientos externos. Ninguna de esas capas introduce retroactivamente los
objetos que la construcción afirma generar.

El orden rector es
\[
\begin{aligned}
\text{unidad discreta}
&\longrightarrow
\text{incidencia orientada}
\longrightarrow
\text{transporte con memoria}\\
&\longrightarrow
\text{número}
\longrightarrow
\text{magnitud}.
\end{aligned}
\]
Cada flecha tiene dominio, codominio y datos conservados. La palabra
\emph{origen} se utiliza sólo cuando esa dependencia se exhibe; una semejanza,
una proximidad decimal o una igualdad cardinal se describen con un estatuto
menor.

\begin{thesisbox}
HMT--MD construye una genealogía matemática conjunta de número, memoria y
magnitud. Los números y constantes se obtienen internamente cuando así lo
establecen sus teoremas; magnitudes, campos y partículas aparecen, según el
caso, como caracteres internos, transducciones, interfaces físicas o
programas de realización. Las teorías físicas conocidas proporcionan después
cartas de realización y contraste, preservando la dirección causal de la
construcción. En la rama excepcional, APP contiene el esqueleto
\(A_2^{12}\), TPK transporta el mismo \(C_3\) sobre caminos, el acoplamiento
orientado APP--Fock produce \(54\), y el alineamiento APP--Witt conduce al
vecino de Leech, \(V^\natural_{\mathrm{HMT}}\), el Monstruo y Moonshine.
\end{thesisbox}

La forma causal completa del cierre excepcional parte del TPK en su perfil
enriquecido y conserva la bifurcación del código \(C_W\):
\[
\mathrm{APP}\to\mathrm{TRIT}\to\mathrm{TPK}_{\rm enr}
\longrightarrow
\begin{cases}
\text{Fock orientado}\longrightarrow54=[q]t_3,\\
 C_W\longrightarrow W_{12}\longrightarrow W_{24},\\
 C_W\xrightarrow{\phi}N(A_2^{12})
 \xrightarrow{v_\alpha}\Lambda_{24}
 \longrightarrow V_\Lambda\longrightarrow V^\natural
 \longrightarrow\mathbb M\longrightarrow J.
\end{cases}
\]
El diagrama es ramificado porque conserva los tipos de sus objetos. Su
compatibilidad se demuestra después de aplicar las representaciones,
transportes y trazas declarados; las dos ramas no se sustituyen ni se
presentan como una sola flecha lineal.

Esta prioridad causal no homogeneiza la fuerza probatoria de todas las
salidas. La dinámica finita APP--TRIT--TPK construye coinductivamente una
sección infinita común cuyas evaluaciones cilíndricas son
\(\pi\), \(e\) y \(\varphi\). Semillas, lectores, orientación, reglas de
prolongación y datos de frontera quedan declarados antes de la generación;
los controles decimales se abren después. Las ejecuciones a una precisión
finita son truncamientos reproducibles de esa sección, no el alcance de la
ley. Las interfaces físicas posteriores conservan, por separado, la fuerza
que corresponda a cada entrelazador.

\section*{Una historia de cambios de estatuto}

La estructura narrativa puede condensarse en la sucesión
\[
u_{\text{objeto}}
\longrightarrow
u_{\text{cuenta}}
\longrightarrow
u_{\text{magnitud}}
\longrightarrow
u_{\text{razón}}
\longrightarrow
\pm u
\longrightarrow
iu
\longrightarrow
u_{\mathbb R}
\longrightarrow
u_{\text{acción}}
\longrightarrow
u_{\text{metrológica}}.
\]
No se propone como una cronología literal y única. Es una arquitectura
filosófica para leer una historia en la que distintas culturas y disciplinas
construyeron operaciones parcialmente independientes. El aporte diferencial
de HMT consiste en reunir posición, orientación, operación y memoria sin
identificarlas prematuramente.

\begin{readingbox}
La matemática aprendió sucesivamente a contar la unidad, medirla, orientarla,
hacerla girar, completarla y utilizarla como unidad física. HMT construye
conjuntamente esas operaciones y conserva la memoria de cada transición.
\end{readingbox}

\section*{Doble registro de exposición: formal y genealógico}

La obra utiliza deliberadamente dos registros. El registro formal especifica
conjuntos, categorías, módulos, cociclos, operadores, límites y
representaciones. El registro narrativo explica por qué esos objetos responden
a una pregunta histórica y qué información recuperan. Ninguno sustituye al
otro. Una fórmula sin genealogía puede ser correcta y conceptualmente opaca;
una metáfora sin morfismo puede ser fértil y matemáticamente indeterminada.

Cuando se afirma que APP es un mapa, TRIT un compás y TPK una ley de
transporte, la metáfora sólo orienta. Las definiciones exactas son,
respectivamente, una célula aritmética provista de caminos y dos hojas, una
reducción ternaria orientada y una categoría de prolongaciones enriquecidas.
Cuando se habla de materia como espacio interiorizado o de gravedad como
transversalidad, la frase conserva estatuto interpretativo. Sólo se promoverá
a realización física cuando se construya el morfismo correspondiente y se
demuestren sus propiedades.

\section*{La prueba como genealogía reproducible}

Cada bloque central se presenta en este orden:

\begin{enumerate}
\item objeto y dominio;
\item regla o mapa;
\item derivación;
\item premisas;
\item estatuto de procedencia y fuerza probatoria;
\item control negativo o falsador;
\item relación con la cadena general.
\end{enumerate}

La exigencia no responde a prudencia retórica. Impide tres regresiones
matemáticas: confundir una salida visible con el estado enriquecido que la
produce; usar una magnitud conocida para seleccionar el lector que después
pretende derivarla; identificar objetos sólo porque comparten cardinal.

\section*{Simetría del estándar comparativo}

Toda comparación externa se somete a una regla simétrica. No se exigirá a
HMT--MD que derive unidades, una ontología física o un selector experimental
como condición para reconocer un teorema interno si el marco comparado adopta
esas mismas entidades como primitivas, convenciones o datos ajustados.
Recíprocamente, la ausencia de una derivación análoga en el marco comparado no
constituye evidencia a favor de HMT--MD. Cada construcción se juzga en su
dominio declarado; toda identificación transversal debe aportar tipos,
mapas, hipótesis, independencia respecto del blanco y falsadores comparables.
Esta simetría impide tanto la penalización por originalidad como la concesión
de privilegios probatorios.

\section*{Lo que significa una obra autocontenida}

El cuerpo incluye las definiciones y demostraciones necesarias para seguir la
cadena publicada. Los programas y certificados constituyen testigos
reproducibles; las notas de procedencia permiten localizar las fuentes, pero
no sustituyen el argumento. La edición se trabaja en un árbol independiente y
conserva, como testigo sellado, el volumen del que procede. Toda reorganización
debe mantener la correspondencia entre afirmación, prueba, certificado y
estatuto.
